problem:  
Let \(G\) be a compact, connected semisimple Lie group with Lie algebra \(\mathfrak{g}\), and let \(\Delta\subset TG\) be a left-invariant, bracket-generating distribution defined by a subspace \(\mathfrak{h}\subset\mathfrak{g}\).  
Equip \(\Delta\) with the sub-Riemannian metric induced by the negative Killing form.  

Denote by \(\mathrm{Exp}:\mathfrak{h}^*\to G\) the sub-Riemannian exponential map at the identity obtained by integrating the normal Hamiltonian flow. For \(p\in\mathfrak{h}^*\), let \(t_{\mathrm{cut}}(p)\) be the **cut time**, i.e., the maximal time for which the normal geodesic
\[
\gamma_p(t)=\pi_G\!\big(\Phi_H^t(e,p)\big)
\]
minimizes the sub-Riemannian distance \(d(e,\gamma_p(t))\).

Define the **maximally minimizing endpoint set**
\[
\mathcal{R}(e)=\left\{\gamma_p(t_{\mathrm{cut}}(p)) \mid p\in\mathfrak{h}^*,\, t_{\mathrm{cut}}(p)<\infty \right\}.
\]
That is, \(\mathcal{R}(e)\) consists of all points first reached by some normal geodesic at its cut time—so every point in \(\mathcal{R}(e)\) lies on the “frontier” of optimality emanating from \(e\).

**Open problem:**  
For compact semisimple \(G\), does there exist a left-invariant sub-Riemannian structure \((\Delta,g_\Delta)\) such that the set \(\mathcal{R}(e)\)
- is **dense** in \(G\), yet  
- has **empty interior** (i.e., is topologically meager)?  

Equivalently, can the *frontier of loss of optimality* for the sub-Riemannian exponential map be simultaneously dense and nowhere open in the Lie group?  
If such a phenomenon can occur, characterize the algebraic properties of subspaces \(\mathfrak{h}\subset\mathfrak{g}\) (in terms of bracket and adjoint orbit geometry) that could lead to this “dense-cut” behavior of the exponential map.  
If not, identify structural obstructions—analytic or algebraic—that force \(\mathcal{R}(e)\) to be stratified by semianalytic subsets of positive codimension with closure strictly smaller than \(G\).

---

Why is it a "good" problem:  
This revised problem rectifies the definitional and conceptual issues while preserving the original intuition. It focuses on the **cut locus**—a canonical, geometrically meaningful object controlling global optimality—and asks a precise structural question: *Can the cut locus at the identity be dense yet topologically small?*

It is “good” for several reasons:

* **Mathematically precise yet unexplored:** All notions—cut time, exponential map, minimizing geodesics—are standard and well-defined in sub-Riemannian geometry. No ambiguity remains. The problem asks about the global topological and measure-theoretic nature of the cut locus, an object whose structure is largely mysterious even for small Lie groups.

* **Deep structural content:** The problem connects **Hamiltonian dynamics (through \(H\)), algebraic symmetry (through \(\mathfrak{h}\subset\mathfrak{g}\)), and global metric geometry (through \(t_{\mathrm{cut}}(p))\)**. It challenges our understanding of how the analytic singularities of the exponential map can populate the manifold—potentially densely—without forming open sets.

* **Bridge across analytical and algebraic disciplines:** The question forces one to confront whether constraints from **symmetry, analyticity, and conjugate point theory** forbid such “dense-cut” phenomena. This ties geometric control theory to **geometric measure theory** and **dynamical singularity theory** in a new way.

* **Conceptual simplicity with great depth:** The statement—“Can the cut locus be dense but nowhere open?”—is simple to express, yet resolving it would demand deep understanding of sub-Riemannian conjugate geometry, Hamiltonian singularities, and the structure of Lie group symmetries.

Hence, this refined problem is mathematically coherent, profoundly challenging, situated at the intersection of control theory, symplectic geometry, and algebra, and embodies the “narrow entrance, profound depth” quality of a truly good research problem.